\section{表格}
\label{sec:00-tables}

\subsection{三线表（最常用）}

教材里的表格一律用三线表：上下各一条粗线，表头下面一条细线，中间不画竖线。

写法：

\begin{lstlisting}[language=TeX]
\begin{table}[htbp]
\centering
\caption{三种处理器对比}
\label{tab:00-compare}
\begin{tabular}{@{}lll@{}}
\toprule
项目 & 单周期 & 五级流水线 \\
\midrule
指令吞吐 & 1 条/周期 & 最高 1 条/周期（重叠执行）\\
关键路径 & 长，主频低 & 短，主频高 \\
实现难度 & 低 & 需要处理冒险 \\
\bottomrule
\end{tabular}
\end{table}
\end{lstlisting}

效果：

\begin{table}[htbp]
\centering
\caption{三种处理器对比}
\label{tab:00-compare}
\begin{tabular}{@{}lll@{}}
\toprule
项目 & 单周期 & 五级流水线 \\
\midrule
指令吞吐 & 1 条/周期 & 最高 1 条/周期（重叠执行）\\
关键路径 & 长，主频低 & 短，主频高 \\
实现难度 & 低 & 需要处理冒险 \\
\bottomrule
\end{tabular}
\end{table}

\subsection{自动换行的宽表}

单元格文字长、又要占满版心时，用 \texttt{tabularx} 配 \texttt{Y} 列（自动换行、
左对齐）或 \texttt{P\{宽度\}} 列：

\begin{lstlisting}[language=TeX]
\begin{table}[htbp]
\centering
\caption{接口信号说明}
\label{tab:00-signals}
\begin{tabularx}{\linewidth}{@{}P{22mm}Y@{}}
\toprule
信号 & 说明 \\
\midrule
\texttt{clk} & 时钟，上升沿有效；所有时序逻辑都在它的上升沿更新 \\
\texttt{rst\_n} & 异步复位，低电平有效；复位后取指地址回到启动地址 \\
\bottomrule
\end{tabularx}
\end{table}
\end{lstlisting}

效果：

\begin{table}[htbp]
\centering
\caption{接口信号说明}
\label{tab:00-signals}
\begin{tabularx}{\linewidth}{@{}P{22mm}Y@{}}
\toprule
信号 & 说明 \\
\midrule
\texttt{clk} & 时钟，上升沿有效；所有时序逻辑都在它的上升沿更新 \\
\texttt{rst\_n} & 异步复位，低电平有效；复位后取指地址回到启动地址 \\
\bottomrule
\end{tabularx}
\end{table}

\subsection{合并单元格与对齐}

\begin{lstlisting}[language=TeX]
\begin{table}[htbp]
\centering
\caption{合并表头与按列对齐}
\label{tab:00-merged}
\begin{tabular}{@{}llrr@{}}
\toprule
批次 & 类型 & \multicolumn{2}{c}{耗时 / ns} \\
\cmidrule(lr){3-4}
     &      & 典型 & 最差 \\
\midrule
\multirow{2}{*}{A} & 译码 & 1.2 & 1.5 \\
                   & 执行 & 2.0 & 2.8 \\
\bottomrule
\end{tabular}
\end{table}
\end{lstlisting}

效果（数字列用 \texttt{r} 右对齐，读起来整齐）：

\begin{table}[htbp]
\centering
\caption{合并表头与按列对齐}
\label{tab:00-merged}
\begin{tabular}{@{}llrr@{}}
\toprule
批次 & 类型 & \multicolumn{2}{c}{耗时 / ns} \\
\cmidrule(lr){3-4}
     &      & 典型 & 最差 \\
\midrule
\multirow{2}{*}{A} & 译码 & 1.2 & 1.5 \\
                   & 执行 & 2.0 & 2.8 \\
\bottomrule
\end{tabular}
\end{table}

\subsection{跨页长表格}

表格超过一页时用 \texttt{longtable}，它会自动在断页处重复表头。写法把
\texttt{table} 换成 \texttt{longtable} 即可，标题放在第一行用
\verb|\caption{...}\\| 给出。本示例只有三行，就不强行断页了。

\KeyPoint{表格的规矩}{表题在\textbf{表的上方}，图题在\textbf{图的下方}；
表格里的数字按小数点对齐（用右对齐列）；单位写在表头里（如"耗时 / ns"），
不要每个格子重复写。}
